
\paragraph{\textbf{Corollary 2.8.}}
\SGARefTarget{sga3:xiv:2.8}
\emph{Under the hypotheses of \SGARefLink{sga3:xiv:2.4}{2.4}, let \(\mathfrak d\) be a Cartan
subalgebra of \(\mathfrak g\). For every \(s\in S\) such that
\(k(s)\) is infinite, there exist an open neighborhood \(V\) of
\(s\) and a regular section \(a\) of \(\mathfrak g\) over \(V\)
such that
\[
  \mathfrak d|_V=\Nil\!\left(a,\mathfrak g|_V\right);
\]
equivalently, \(a\) is a section of \(\mathfrak d|_V\).}

Indeed, since \(k(s)\) is infinite, there is a regular element
\(a_0\) of \(\mathfrak g(s)\) contained in \(\mathfrak d(s)\).
After replacing \(S\) by an open neighborhood \(U\) of \(s\), we
may extend \(a_0\) to a section \(a\) of \(\mathfrak d\) over
\(U\). The endomorphism induced by \(\operatorname{ad}(a)\) on the
fiber
\[
  \mathfrak g(s)/\mathfrak d(s)
\]
is an automorphism. Shrinking \(U\) to a smaller open neighborhood
\(V\), we may therefore suppose that
\(\operatorname{ad}(a)_{\mathfrak g/\mathfrak d}\) is an
automorphism. It follows that \(a\) is a quasi-regular section of
\(\mathfrak g|_V\) satisfying the required condition.

Let \(\mathfrak g\) be as in \SGARefLink{sga3:xiv:2.4}{2.4}. Inspection of its Killing
polynomial immediately shows that the function
\[
  s\longmapsto\text{the nilpotent rank of }\mathfrak g(s)
\]
on \(S\) is upper semicontinuous. We shall be especially interested
in the case in which this function is continuous, that is, locally
constant. Here are several variants of this property.
